\chapter{Формирование требований к системе}

\section{Анализ существующих подходов}

Современные компьютерные сети подвержены широкому спектру угроз, включая DDoS-атаки, сканирование портов, распространение вредоносного ПО и утечки данных. Для защиты от этих угроз применяются системы обнаружения вторжений (IDS), которые делятся на три основные категории.

\subsection{Сигнатурные методы}

Сигнатурные методы сравнивают сетевой трафик с базой известных сигнатур атак. К наиболее известным инструментам относятся:

\begin{itemize}
    \item \textbf{Snort} — популярная OpenSource IDS, использующая правила для обнаружения известных атак. Поддерживает анализ протоколов, поиск по содержимому пакетов и обнаружение аномалий.
    \item \textbf{Suricata} — многопоточная IDS/IPS с поддержкой HTTP, TLS и других протоколов, обеспечивающая высокую производительность за счёт распараллеливания обработки.
    \item \textbf{Zeek (Bro)} — платформа анализа сетевого трафика, фокусирующаяся на семантическом анализе и генерации структурированных логов.
\end{itemize}

Преимуществом сигнатурных методов является низкий уровень ложных срабатываний. Основной недостаток — невозможность обнаружения новых, ранее неизвестных атак.

\subsection{Статистические методы}

Статистические методы выявляют аномалии путём анализа отклонений от нормального профиля сетевой активности:

\begin{itemize}
    \item \textbf{Пороговые методы} (threshold-based) — обнаружение превышения нормальных показателей трафика (количество пакетов, объём данных).
    \item \textbf{Анализ временных рядов} — выявление сезонности и трендов сетевой активности с использованием скользящего среднего и авторегрессии.
    \item \textbf{Методы скользящего окна} — вычисление статистических характеристик (среднее, дисперсия, стандартное отклонение) в скользящих временных интервалах.
\end{itemize}

\subsection{Методы машинного обучения}

Методы машинного обучения являются наиболее перспективным направлением для обнаружения аномалий сетевой активности.

\textbf{Isolation Forest} (Изолирующий лес) — метод, основанный на том, что аномалии легче изолировать случайными разбиениями, чем нормальные объекты. Алгоритм строит ансамбль случайных деревьев решений; чем меньше разбиений требуется для изоляции точки, тем выше вероятность, что она является аномалией. Преимущества: эффективность для многомерных данных, устойчивость к выбросам, линейная сложность.

\textbf{LOF (Local Outlier Factor)} — метод, основанный на локальной плотности объектов. Для каждой точки вычисляется отношение её локальной плотности к средней плотности k ближайших соседей. Точки с аномально низкой плотностью (LOF значительно больше 1) считаются выбросами. Преимущества: не требует предположений о распределении данных, выявляет локальные аномалии.

\textbf{One-Class SVM} — метод опорных векторов для одного класса. Строит разделяющую гиперплоскость, максимально отделяющую нормальные данные от начала координат в пространстве признаков, преобразованном с помощью ядер. Преимущества: эффективен для данных с чёткой границей нормального класса.

\begin{table}[H]
\centering
\caption{Сравнительный анализ подходов к обнаружению аномалий}
\label{tab:comparison}
\begin{tabular}{p{4cm}p{5.5cm}p{5.5cm}}
\toprule
\textbf{Метод} & \textbf{Преимущества} & \textbf{Недостатки} \\
\midrule
Сигнатурный & Низкий уровень ложных срабатываний; простота настройки; высокая скорость обработки & Не обнаруживает новые атаки; требует постоянного обновления базы сигнатур \\
\addlinespace
Статистический & Простота реализации; интерпретируемость результатов; низкие требования к ресурсам & Чувствителен к выбору порогов; сложность настройки для гетерогенных сетей \\
\addlinespace
Isolation Forest & Эффективен для многомерных данных; устойчив к выбросам; линейная сложность & Требует настройки гиперпараметров (n\_estimators, contamination) \\
\addlinespace
LOF & Не требует предположений о распределении; выявляет локальные аномалии & Чувствителен к выбору k-соседей; высокая вычислительная сложность \\
\addlinespace
One-Class SVM & Хорошо работает на размеченных данных; эффективен для чётких границ класса & Медленен на больших объёмах; чувствителен к выбору ядра \\
\bottomrule
\end{tabular}
\end{table}

\subsection{Потоковая обработка данных}

Для обработки сетевого трафика в реальном времени используются платформы потоковой обработки:

\begin{itemize}
    \item \textbf{Apache Kafka} — распределённая платформа для сбора и передачи потоковых данных с высокой пропускной способностью.
    \item \textbf{Apache Spark Streaming} — фреймворк для потоковой обработки, основанный на микро-пакетной архитектуре.
    \item \textbf{Apache Flink} — фреймворк для обработки событий в реальном времени с поддержкой exactly-once семантики.
    \item \textbf{Pandas + NumPy} — стек для анализа данных в оперативной памяти, оптимальный для небольших и средних объёмов данных.
\end{itemize}

Для данного проекта выбран стек Python + Pandas + NumPy + Scikit-learn как оптимальный по соотношению функциональности, производительности и простоты разработки.

\section{Функциональные требования}

На основе проведённого анализа сформулированы следующие функциональные требования к разрабатываемой платформе:

\begin{enumerate}[label={\bfseries RF-\arabic*:}]
    \item \textbf{Сбор сетевого трафика} — возможность захвата пакетов в режиме реального времени (live capture) и загрузки данных из PCAP-файлов.
    \item \textbf{Извлечение признаков} — автоматическое извлечение характеристик из сетевых пакетов: IP-адреса источника и назначения, порты, протоколы, размеры пакетов, временные метки, TCP-флаги, TTL.
    \item \textbf{Потоковая агрегация} — группировка пакетов во временные окна с вычислением статистических признаков: количество пакетов, средний размер, стандартное отклонение, энтропия, количество уникальных адресов.
    \item \textbf{Обнаружение аномалий} — реализация алгоритмов машинного обучения (Isolation Forest, LOF, One-Class SVM) и эвристических правил для выявления аномалий.
    \item \textbf{Визуализация результатов} — построение графиков временных рядов с отмеченными аномалиями, confusion matrix, ROC-кривых, распределения оценок аномалий.
    \item \textbf{Модульная архитектура} — возможность независимой замены или добавления новых алгоритмов обнаружения без модификации существующего кода.
    \item \textbf{Сохранение моделей} — возможность сохранения обученных моделей на диск и их последующей загрузки.
    \item \textbf{Генерация тестовых данных} — встроенный генератор синтетического трафика для тестирования, поддерживающий несколько типов атак.
\end{enumerate}

\section{Нефункциональные требования}

\begin{enumerate}[label={\bfseries RNF-\arabic*:}]
    \item \textbf{Производительность} — обработка не менее 1000 пакетов в секунду на стандартном оборудовании.
    \item \textbf{Масштабируемость} — модульная архитектура, позволяющая расширять функциональность без переработки существующего кода.
    \item \textbf{Кроссплатформенность} — поддержка операционных систем Windows, Linux, macOS.
    \item \textbf{Воспроизводимость} — использование детерминированных алгоритмов для обеспечения повторяемости результатов.
    \item \textbf{Качество кода} — соответствие стандартам PEP 8, наличие комментариев и документации.
\end{enumerate}

\section{Выбор технологического стека}

На основе анализа требований и существующих решений был выбран следующий технологический стек:

\begin{table}[H]
\centering
\caption{Выбор технологического стека}
\label{tab:tech}
\begin{tabular}{p{3.5cm}p{3cm}p{8cm}}
\toprule
\textbf{Компонент} & \textbf{Технология} & \textbf{Обоснование} \\
\midrule
Захват трафика & Scapy & Гибкость, поддержка множества протоколов, работа с PCAP \\
Обработка данных & Pandas + NumPy & Мощные структуры данных, богатая экосистема \\
ML-модели & Scikit-learn & Широкий выбор алгоритмов, простота использования \\
Визуализация & Matplotlib + Seaborn & Качественная 2D-графика, настраиваемые графики \\
Хранение результатов & CSV / joblib & Простота, совместимость, быстрая сериализация \\
Генерация данных & Встроенный модуль & Независимость от внешних источников данных \\
\bottomrule
\end{tabular}
\end{table}

\subsection{Обоснование выбора Python}

Выбор языка Python обусловлен несколькими факторами:
\begin{enumerate}
    \item Наличие богатой экосистемы библиотек для анализа данных и машинного обучения (Pandas, NumPy, Scikit-learn).
    \item Простота разработки и читаемость кода, что важно для учебного проекта.
    \item Кроссплатформенность — код работает на Windows, Linux и macOS без модификаций.
    \item Активное сообщество и обширная документация.
\end{enumerate}
